\section*{Limitations}
\label{app:limitations}

\noindent\textbf{Scope of our task formulation}. As mentioned in \S\ref{sec:intro}, we formulate the implicit reasoning problem as induction and application of inference rules from a mixture of atomic and inferred facts. This may not apply to the full spectrum of reasoning which has a range of different types and meanings~\citep{huang-chang-2023-towards}. Still, our formulation could capture the nature of a wide range of reasoning problems, and crucially, we believe that it is a good conceptualization of certain aspects of language model (pre-)training, where the model needs to both memorize the ``atomic'' world knowledge and also induce generalizable rules from the massive amount of records of human activities, which allow the model to connect and reason over knowledge, and ultimately help with humans.

\noindent\textbf{Abstract nature and connection with practice}. Our study has an abstract nature, which is required for the complete solidity and rigor of our experiments and evaluation. We also strive to make sure that the results are robust to different setups closer to practice through additional experiments (Appendix~\ref{app:scaling},\ref{app:tokenization},\ref{app:other_results}). Still, there are certain distances from our settings to those in practice. However, we believe that it is far more important to build solid understandings, even having distances with practice, than to draw conclusions or make claims that are closer to practice but questionable due to insufficient control over data and evaluations.